\makeabstract
{
Effect of Ethyl-3-Hydroxybutyrate on Viability of Pro-inflammatory Macrophages for Sepsis Treatment
}
{
Benjamin Edwards (1), Gracy Rosario, PhD. (2), Gelilla Daniel, MD (3), Cuthbert Simpkins, MD (2,3)
}
{
\textsuperscript{1}Amherst College, Amherst, MA. \\
\textsuperscript{2}University of Missouri-Kansas City School of Medicine, Kansas City, MO. \\
\textsuperscript{3}Vivacelle Bio, Inc. Kansas City, MO
}
{
Sepsis is an inflammatory condition, caused by aberrant host immune response to microbial infection, leading to increased tissue damage and mortality. The increased incidence of antibiotic resistance worldwide and lack of candidate drugs for sepsis treatment has increased mortality and morbidity in the USA \& worldwide. Vivacelle Bio has been involved in developing phospholipid nanoparticles for treatment of sepsis, of which VBI-S has cleared Phase IIa clinical trials. Recently, we have also found a ketone body, ethyl-3-hydroxybutyrate (EHB), that can increase survival of severe sepsis in a cecal ligation and rat laceration model. In this model, CLL rats treated with 125mM EHB in aqueous phase at 8\% of the rats body weight led to complete recovery within 36h. EHB (20mM) has been shown to demonstrate potent immunomodulatory effects, as it reduces pro-inflammatory cytokine/chemokine gene expression in M1 macrophages. However, it is unknown whether the reduced gene expression in M1 macrophages is linked to cell cytotoxicity by EHB.
}
{
THP-1 cells (50000) were seeded into 96 well plates and differentiated into M1 macrophages. Cells were treated in their respective groups for 24 hours. The post culture supernatant was removed, stained with trypan blue stain, and counted with a hemocytometer to determine the number of dead cells. The adhered cells were terminated with 4N HCl in isopropanol and the A570 was read. Statistical significance was determined by using Brown Forsythe and Welch ANOVA tests. Post Hoc analysis by Mann Whitney U tests were used to examine lower concentrations of EHB.
}
{
It was found that EHB is cytotoxic at 67.5mM concentration, causing both decreased metabollic activity and an increased number of dead cells. However, EHB promotes metabolic activity at 10mM and 20mM concentrations, while also decreasing the number of dead cells, improving overall cell viability. 
}
{
The decrease in pro-inflammatory cytokine/chemokine expression by EHB is caused by gene expression regulation at 10 and 20 mM instead of a cytotoxic effect lowering M1 macrophage populations. These findings suggest that mitigating the dysregulated inflammatory response with EHB could play a critical role in the effective treatment of sepsis. Future studies are necessary to develop a better understanding of the mechanism(s) by which EHB reduces inflammatory responses in M1 macrophages and enhances survival.
}

\newpage
